%==================================================================
% BAB 5 PENUTUP
%  5.1 Kesimpulan
%  5.2 Saran
%==================================================================
\chapter{PENUTUP}
\label{chap:penutup}

%------------------------------------------------------------------
\section{Kesimpulan}
%------------------------------------------------------------------

Berdasarkan hasil perancangan, pengujian, dan pembahasan yang telah
diuraikan, dapat disimpulkan sebagai berikut.
\begin{packed_enum}
  \item Sistem deteksi kesegaran ikan berhasil dirancang bangun
        menggunakan kamera digital di dalam kotak akusisi dengan dua
        lampu LED yang menghasilkan pencahayaan seragam pada seluruh
        citra.
  \item Ekstraksi ciri warna pada ruang HSV menghasilkan enam ciri
        berupa nilai rata-rata dan simpangan baku kanal $H$, $S$, dan
        $V$ yang mampu membedakan tiga tingkat kesegaran ikan.
  \item Jaringan saraf tiruan dengan arsitektur 6-16-8-3 mencapai
        akurasi 92,33\% pada pengujian \textit{5-fold cross validation},
        dengan presisi tertinggi pada kelas segar sebesar 96,94\%.
  \item Kesalahan klasifikasi terbanyak terjadi antara kelas cukup
        segar dan tidak segar karena kemiripan warna insang pada kedua
        kelas tersebut.
\end{packed_enum}

%------------------------------------------------------------------
\section{Saran}
%------------------------------------------------------------------

Beberapa saran yang dapat diberikan untuk pengembangan selanjutnya
adalah sebagai berikut.
\begin{packed_enum}
  \item Penelitian selanjutnya dapat menambahkan ciri tekstur insang
        dan bentuk mata untuk memperkecil kesalahan antara kelas cukup
        segar dan tidak segar.
  \item Sistem dapat dikembangkan menjadi aplikasi telepon pintar agar
        lebih mudah digunakan oleh pedagang di pasar.
  \item Penambahan jumlah dan variasi jenis ikan pada data latih perlu
        dilakukan agar model dapat digunakan secara lebih umum.
\end{packed_enum}
